% Root construction; independent birational-agent audit passed, with the
% coefficient-prime ramification correction incorporated below.
\subsection{Two coefficient primes with the same local elliptic realizations}

\begin{statement}{B.16}{Proved}
{An infinitely ramified pair with common integral local elliptic curves}
There are continuous representations
\[
 \rho_7:G_\Q\longrightarrow\mathrm{GL}_2(\Z_7),\qquad
 \rho_{11}:G_\Q\longrightarrow\mathrm{GL}_2(\Z_{11})
 \tag{B.16.1}
\]
with determinants $\chi_7,\chi_{11}$ and surjective residual
representations, and elliptic curves $E_q/\Q_q$ for every finite
prime $q$, with integral Galois-module isomorphisms
\[
 \rho_\ell|_{G_{\Q_q}}\simeq T_\ell(E_q)
       \qquad(\ell=7,11).
 \tag{B.16.2}
\]
The same curve $E_q$ occurs on both sides of the coefficient choice.
The two ramification sets agree away from $\{7,11\}$ and are
infinite, each of density zero.
At every common unramified prime $q\ne7,11$ their common Frobenius
polynomial is $X^2-a_qX+q$, with $a_q\in\Z$ and
$|a_q|\le2\sqrt q$.
For $q\in\{2,3,7,11\}$ one can take
$E_q:y^2=x(x-1)(x-3)$.  At every newly introduced ramified prime
one can take the literal split Tate curve with parameter $q^{M_q}$,
where $M_q$ is a positive integer specified during the construction.
\end{statement}
\begin{proof}
We prove the common auxiliary-prime and integral-lattice assertions
needed to strengthen the one-coefficient construction B.8.

\emph{The two residual representations.}
Put $E:y^2=x(x-1)(x-3)$ and start with $\bar\rho_\ell=E[\ell]$.
The computation in B.8 gives $\Delta=576$, a transvection from
inertia at $3$, and $a_5=2$.  The discriminant $-16$ of the
Frobenius polynomial at $5$ is a nonsquare modulo both $7$ and
$11$: its residues are $5$ and $6$, respectively.
The transvection and an element moving its fixed line generate
$\mathrm{SL}_2(\F_\ell)$ by the elementary-matrix argument in B.8.
The cyclotomic determinant is surjective.  Thus both residual
images are $\mathrm{GL}_2(\F_\ell)$.
All lifts below have the indicated fixed cyclotomic determinant.

\emph{The exact disjointness used at every auxiliary-prime choice.}
Fix any finite stage of the deformation construction at $\ell=7$
or $11$.  Form the compositum $L_\ell$ of the field cut out by
its finite-level representation and the finitely many fields
cut out by adjoint or dual-adjoint cohomology classes used at
that stage.  The cocycle map defining one such field is the
actual homomorphism
\[
 G_\Q\longrightarrow M\rtimes\mathrm{GL}_2(\F_\ell),
 \qquad g\longmapsto(f(g),\bar\rho_\ell(g)),
 \tag{B.16.3}
\]
with $M=\operatorname{Ad}^0\bar\rho_\ell$ or
$M=\operatorname{Ad}^0\bar\rho_\ell(1)$.
The finite congruence kernel and every kernel in (B.16.3) have
$\ell$-power order.  Consequently there is an exact sequence
\[
 1\longrightarrow P_\ell\longrightarrow
 \operatorname{Gal}(L_\ell/\Q)
 \longrightarrow\mathrm{GL}_2(\F_\ell)\longrightarrow1,
 \qquad P_\ell\text{ an }\ell\text{-group}.
 \tag{B.16.4}
\]
The cyclotomic fields needed for these actions already occur
through the determinant and have the same property.

Now $7\nmid|\mathrm{GL}_2(\F_{11})|=13200$ and
$11\nmid|\mathrm{GL}_2(\F_7)|=2016$.
The Galois group of $L_7\cap L_{11}$ is a quotient of both
groups in (B.16.4).  Its order is prime to $7$ and $11$, so
both kernels $P_\ell$ die in that quotient.
Every elementary unipotent in $\mathrm{SL}_2(\F_\ell)$ has
order $\ell$, and these elements generate that group.
Their images therefore also die.  The common quotient factors
through the determinants on both sides.  Thus
\[
 L_7\cap L_{11}\subset\Q(\mu_7)\cap\Q(\mu_{11})=\Q.
 \tag{B.16.5}
\]
For the last equality, the intersection has degree dividing
$\gcd(6,10)=2$, and the unique quadratic subfields
$\Q(\sqrt{-7})$ and $\Q(\sqrt{-11})$ are different.
Restriction hence gives the actual isomorphism
\[
 \operatorname{Gal}(L_7L_{11}/\Q)
   \xrightarrow{\ \sim\ }
 \operatorname{Gal}(L_7/\Q)\times
 \operatorname{Gal}(L_{11}/\Q).
 \tag{B.16.6}
\]
This proof applies again after adjoining every additional
cohomology field used in a subsequent finite step.

We use the deformation steps of
\cite[Fact~6, Lemmas~7--9, Propositions~10--11 and \S2.2]{BIR_KLR2005}
with their actual local conditions.  At a prime called nice there,
the residual Frobenius has eigenvalue ratio $q$, with
$q\not\equiv\pm1\pmod\ell$.  We select its split-sign component,
whose Frobenius is conjugate to $\operatorname{diag}(q,1)$.
Here is an exact sign-changing element preserving those
cohomology restrictions.  In (B.16.4), the inverse image of
$\{I,-I\}$ is an extension of $C_2$ by an $\ell$-group.
The finite coprime-complement theorem supplies an element $z$
of order two above $-I$.  Its finite-level matrix is literally
$-I$: the idempotent $(I+\rho(z))/2$ reduces to zero modulo
$\ell$, hence is zero in the finite matrix ring.
On both adjoint modules in (B.16.3) its action is trivial.
For every cocycle used there,
$0=f(z^2)=f(z)+z f(z)=2f(z)$, so $f(z)=0$.
Multiplication of the chosen Frobenius by $z$ therefore
switches the two signs while preserving its determinant,
cyclotomic residue and all its cohomology coordinates.
This proves the split-sign refinement of Fact~6 with its
actual maps.  Equations (B.16.4)--(B.16.6) combine every such
choice with niceness and the split sign for the other
coefficient at any fixed current depth.

This additional fixed Chebotarev condition also survives the
two-prime step in Proposition~11 of that source.
Indeed its proof first restricts to a positive-density
Chebotarev set, then chooses a positive-upper-density subset
on which two finite local values are constant.  The further
fields cut out by its classes $f_t,\phi_t$ are still fields
of the form (B.16.4).  They remain disjoint from the fixed
other-coefficient field.  Each independent complementary
condition in the proof therefore removes the same fixed
positive fraction within the restricted Chebotarev set.
The final bound is its density times $\gamma^h$, with
$0<\gamma<1$ and $h\to\infty$, exactly as in that proof.
Thus one or two auxiliary primes still suffice for that step,
with both coefficient conditions retained.

\emph{Local targets and the alternating induction.}
Use the original curve $E$ at $2,3,7,11$.
Enlarge the initial finite ramification set to kill the two
adjoint and dual-adjoint localization kernels, using Lemma~7
of the cited source with the simultaneous nice-prime condition.
At each extra prime prescribe a split Tate curve $q^M$ with
$M$ divisible by $7^2 11^2$; its residual classes agree with
both prescribed split-sign Frobenius classes.
Start from the finite-level lifts supplied by $E$.
The local corrections of Lemma~9 impose these prescribed
characteristic-zero local targets.  Enlarging the set thereafter
keeps the localization kernels zero: a class locally zero at
a newly added prime is unramified there and inflates from the
old set, as in the proof of Lemma~7.

Here is the level bookkeeping.  At stage $m\ge2$ there are
possibly uncorrected lifts modulo $7^m$ and $11^m$ of the
already corrected level $m-1$ maps.  There is a finite set
$S$ of frozen common local curves, whose reductions agree
with the corrected level $m-1$ maps.  There is a finite set
$R$, disjoint from $S$, with fixed common integral Frobenius
polynomials agreeing at level $m-1$.
First apply the level-$m$ correction and two-prime step at $7$.
Every new auxiliary prime is also chosen nice with split sign
for the current, uncorrected representation modulo $11^m$.
At the end of this correction the new $7$-adic local matrix has
the form, modulo $7^m$,
\[
 \rho_{7,m}(\sigma_q)=\begin{pmatrix}q&0\\0&1\end{pmatrix},
 \qquad
 \rho_{7,m}(\tau_q)=\begin{pmatrix}1&B_7\\0&1\end{pmatrix}.
 \tag{B.16.7}
\]
Choose a positive integer $M_q$ with
$v_7(M_q)=v_7(B_7)$ when $B_7\ne0$ modulo $7^m$, and
$v_7(M_q)\ge m$ otherwise; put $v_{11}(M_q)\ge m$.
Diagonal unit conjugacy makes (B.16.7) the reduction of
$T_7(\operatorname{Tate}(q^{M_q}))$.
Its $11$-adic reduction is unramified and is the selected
matrix $\operatorname{diag}(q,1)$ modulo $11^m$.
Freeze this same curve at $q$ and add $q$ to $S$.

Next carry out the level-$m$ correction at $11$, prescribing
the common curves at this enlarged $S$.
Every new auxiliary prime in this step is selected nice with
split sign for the \emph{corrected} representation modulo $7^m$.
After the $11$ correction choose $M_q$ with the valuation
prescribed by its nonzero $B_{11}$, or at least $m$ when that
residue is zero, and with $v_7(M_q)\ge m$.
It now gives a common Tate curve whose reductions match
both corrected level-$m$ representations.
These choices preserve all previously fixed local curves
and Frobenius polynomials.  The independent field comparison
(B.16.6) justifies every auxiliary-prime request in the two steps.

Both corrected level-$m$ maps have characteristic-zero local
lifts at every prime of the enlarged $S$.
The vanishing adjoint localization kernels show that their
global lifting obstructions to level $m+1$ are zero.
Choose any such lifts unramified outside $S$.
They need not yet agree with the targets modulo $\ell^{m+1}$:
the next stage corrects them by $\ell^m$-valued cochains,
preserving the corrected level $m$.  This is the precise
compatibility used to take the inverse limit.

For infinitude, at every stage also select one fresh simultaneous
nice prime outside all finite sets under consideration, and
prescribe $q^{77^m}$ there before the first correction.
Its finite-level unramified split-sign matrices agree with the
current representations.  It enters $S$ permanently, and its
eventual common Tate parameter is nonzero.  Thus infinitely
many primes actually ramify in both limiting representations.

\emph{The common good-prime polynomials.}
For a concrete exhaustive choice put $C_m=77^{2m}$.
At the first stage prescribe the polynomials of $E$ at all
primes at most $C_2$ outside $S$.
At stage $m>2$, add to $R$ every prime at most $C_m$ not
already in $S\cup R$.  Such a newly added prime satisfies
$q>C_{m-1}$.  The two existing traces specify one residue
class modulo $77^{m-1}$ by the Chinese remainder theorem.
Since
\[
                  2\sqrt q>2\cdot77^{m-1},
 \tag{B.16.8}
\]
there is a nonzero integer $a_q$ in that residue class with
$|a_q|<2\sqrt q$.  Fix its polynomial for both coefficient
primes and impose it by the level correction just described.
All auxiliary primes can be required to exceed $C_m$.
Every rational prime eventually belongs to $S\cup R$,
and the common polynomials on $R$ persist through the
continuous inverse limit.  This gives the two actual maps
(B.16.1), with their common infinite ramification set away
from the two coefficient primes.  At $7$, $\rho_7$ is
ramified through its cyclotomic determinant whereas
$\rho_{11}=T_{11}(E)$ is unramified; the roles are reversed
at $11$.  Thus these two explicitly retained exceptions
are essential in the statement.
Their residual absolute irreducibility makes both limiting
representations semisimple, so
\cite[Theorem~1]{BIR_KhareRajan2001} gives density zero.

\emph{One integral local curve for both coefficient lattices.}
At a new bad prime $q$, both full representations have the
split Tate matrices with inertia entries of the same valuations
as the fixed integer $M_q$.  Their diagonal unit conjugators
rescale the entries to the literal integer $M_q$.
The parameter $q^{M_q}$ therefore realizes both integral
modules; there is no independent unit-parameter choice.
The Tate torsion generators are the compatible roots of unity
and roots of $q^{M_q}$, exactly as in B.8.

At $q\in R$, Waterhouse's theorem as used in B.8 supplies
an elliptic curve over $\F_q$ with polynomial $X^2-a_qX+q$.
Lift it to a good-reduction curve $E'_q/\Q_q$.
For each $\ell=7,11$, its rational Tate module is isomorphic
to the unramified rational representation of $\rho_\ell$:
the polynomial has distinct roots since $q$ is prime.
Inside the rational Tate space let $L_\ell$ be the stable
lattice furnished by $\rho_\ell$.  Multiply it by a scalar
power of $\ell$ so that
\[
 T_\ell(E'_q)\subset L_\ell
                   \subset\ell^{-h_\ell}T_\ell(E'_q).
 \tag{B.16.9}
\]
Then $H_\ell=L_\ell/T_\ell(E'_q)$ is a Galois-stable
finite subgroup of $E'_q[\ell^{h_\ell}]$.
The subgroups $H_7,H_{11}$ have coprime orders.  Their sum
is a finite subgroup defined over $\Q_q$, and the quotient
\[
                      E_q=E'_q/(H_7+H_{11})
 \tag{B.16.10}
\]
has, under the rational isogeny identifications, exactly the
two lattices in (B.16.9).  At one coefficient the other
subgroup has prime-to-$\ell$ order and induces an isomorphism
of integral Tate modules.  Both subgroups extend to finite
etale subgroups of the good integral model, because
$q\notin\{7,11\}$, so the quotient again has good reduction.
This proves (B.16.2), including Frobenius matrices whose
reductions are scalar.  At $2,3,7,11$ the frozen curve was
the stated $E$ from the start.  At the real place both
representations have an involution of determinant $-1$;
since $2$ is a unit in both coefficient rings they are
integrally conjugate to $\operatorname{diag}(1,-1)$,
also realized by $E/\R$.
\end{proof}

\paragraph{Exact comparison with the section problem.}
For a full family with common local curves, the forgetful map is
\[
 \bigl((\rho_\ell)_\ell,(E_q)_q,(\iota_{q,\ell})_{q,\ell}\bigr)
 \longmapsto
 \bigl(\rho_7,\rho_{11},(E_q)_q,
              (\iota_{q,7},\iota_{q,11})_q\bigr).
 \tag{B.16.11}
\]
It preserves continuity, the selected integral local modules,
and their Frobenius polynomials.  The construction proves
infinite ramification for an object in this exact target.
An inverse image under (B.16.11), and a lift through the
Legendre monodromy map from an actual tripod section,
are not constructed here.  Thus the full A4 question retains
its all-coefficient and fundamental-group lifting requirements.

The exact program \path{code/two_coefficient_elliptic_certificate.py}
and its retained output verify the two residual finite-group
calculations, the cross-prime orders, the common Tate-matrix
conjugacies and the CRT trace bound.  They do not simulate
the infinite deformation construction.
